
\textbf{Principaux domaines abordés :} suites, fonctions, primitives

\medskip

%\emph{Cet exercice est un questionnaire à choix multiples.\\
%Pour chacune des questions suivantes, une seule des quatre réponses proposées est exacte.\\
%Une réponse fausse, une réponse multiple ou l'absence de réponse à une question ne rapporte ni n'enlève de point.\\
%Pour répondre, indiquer sur la copie le numéro de la question et la lettre de la réponse choisie.\\
%Aucune justification n'est demandée.}
%
%\medskip

\begin{enumerate}
\item %On considère la suite $\left(u_n\right)$ définie pour tout entier naturel $n$ par

\[u_n = \dfrac{(- 1)^n}{n + 1}.\]

%On peut affirmer que:
%\begin{center}
%\begin{tabularx}{\linewidth}{X X}
%\textbf{a.~~} la suite $\left(u_n\right)$ diverge vers $+\infty$. &\textbf{b.~~} la suite $\left(u_n\right)$ diverge vers $-\infty$.\\
%\textbf{c.~~}  la suite $\left(u_n\right)$ n'a pas de limite. &\textbf{d.~~} la suite $\left(u_n\right)$ converge.
%\end{tabularx}
%\end{center}
On a $- 1 \leqslant (- 1)^n \leqslant 1$, donc puisque $n+1 \geqslant 1 > 0$ et que $\dfrac{1}{n  1} > 0$, on en déduit que pour tout naturel :

\[- \dfrac{1}{n+1} \leqslant u_n \leqslant \dfrac{1}{n+1}\]

La suite $\left(u_n\right)$ encadrée par deux suites ayant pour limite 0 a pour limite 0.
\begin{center} \decosix \decosix \decosix  \end{center}
\end{enumerate}

Dans les questions 2 et 3, on considère deux suites $\left(v_n\right)$  et $\left(w_n\right)$ vérifiant la relation : 

\[w_n = \e^{- 2v_n} + 2.\]

%\begin{enumerate}[resume]
%\item  Soit $a$ un nombre réel strictement positif. On a $v_0 = \ln (a)$.
%
%\begin{center}
%\begin{tabularx}{\linewidth}{X X}
%\textbf{a.~~} $w_0 = \dfrac{1}{a^2}  +2$&\textbf{b.~~}$w_0 = \dfrac{1}{a^2  +2}$\\
%\textbf{c.~~} $w_0 = -2a +2$&\textbf{d.~~}$w_0 = \dfrac{1}{- 2a} + 2$
%\end{tabularx}
%\end{center}
On a $w_0 = \e^{- 2v_0} + 2 = \e^{- 2\ln (a)} + 2 = \dfrac{1}{\e^{2\ln(a)}} + 2=  \dfrac{1}{\e^{\ln \left(a^2\right)}} + 2 = \dfrac{1}{a^2} + 2$.

\begin{enumerate}[resume]
\item On sait que la suite $\left(v_n\right)$ est croissante. On peut affirmer que la suite $\left(w_n\right)$ est :

%\begin{center}
%\begin{tabularx}{\linewidth}{X X}
%\textbf{a.~~}décroissante et majorée par 3.&\textbf{b.~~}décroissante et minorée par 2 .\\
%\textbf{c.~~}croissante et majorée par 3 .&\textbf{d.~~}croissante et minorée par 2.
%\end{tabularx}
%\end{center}
On a successivement :

$v_n \leqslant v_{n+1} \Rightarrow 2v_n \leqslant 2v_{n+1} \Rightarrow - 2v_{n+1} \leqslant - 2v_n\Rightarrow \e^{- 2v_{n+1}} \leqslant \e^{- 2v_n}$ par croissance de la fonction exponentielle et enfin $\e^{- 2v_{n+1}} + 2 \leqslant \e^{- 2v_n} + 2$, soit $w_{n+1} \leqslant w_n$ : la suite $\left(w_n\right)$ est décroissante.

D'autre part on sait que quel que soit le réel$\alpha, \: \e^{\alpha} > 0$ donc  la suite est minorée par 2.
\item %On considère la suite $\left(a_n\right)$ ainsi définie :

\[a_0 = 2 \:\:\text{et, pour tout entier naturel}\: n, \quad a_{n+1} = \dfrac13a_n + \dfrac83.\]

%Pour tout entier naturel $n$,on a :
%\begin{center}
%\begin{tabularx}{\linewidth}{X X}
%\textbf{a.~~}$a_n = 4 \times \left(\dfrac13\right)^n - 2$&\textbf{b.~~}$a_n = - \dfrac{2}{3^n} + 4$\\
%\textbf{c.~~}$a_n = 4 - \left(\dfrac13\right)^n$ & \textbf{d.~~} $a_n = 2 \times \left(\dfrac13\right)^n + \dfrac{8n}{3}$
%\end{tabularx}
%\end{center}
Soit la suite $\left(b_n\right)$ définie pour tout naturel $n$ par $b_n = a_n + \alpha$.

Alors $b_{n+1} = a_{n+1} + \alpha = \dfrac13a_n + \dfrac83 + \alpha = \dfrac13\left(b_n - \alpha \right) + \dfrac83 + \alpha = \dfrac13b_n - \dfrac{\alpha}{3} + \dfrac83 + \alpha = \dfrac13b_n + \dfrac{8 + 2\alpha}{3}$.

Prenons $\alpha = - 4$, alors $b_{n+1} = \dfrac13b_n$ : cette égalité montre que la suite $\left(b_n\right)$ est une suite géométrique de raison $\dfrac13$ et de premier terme $b_0 = a_0 + \alpha = 2 - 4 = - 2$.

On sait qu'alors que, quel que soit le naturel $n$, \: $b_n = b_0 \times \left(\dfrac{1}{3}\right)^n = - 2 \times \left(\dfrac{1}{3}\right)^n$.

Finalement $b_n = a_n - 4 \iff a_n = 4 + b_n = 4 - 2\times \left(\dfrac{1}{3}\right)^n$.
\item On considère une suite $\left(b_n\right)$ telle que, pour tout entier naturel $n$, on a :

\[b_{n+1} = b_n + \ln \left(\dfrac{2}{\left(b_n \right)^2 + 3}\right).\]

%On peut affirmer que :
%
%\begin{center}
%\begin{tabularx}{\linewidth}{X X}
%\textbf{a.~~} la suite $\left(b_n\right)$ est croissante.&\textbf{b.~~}la suite $\left(b_n\right)$ est décroissante.\\
%\textbf{c.~~} la suite $\left(b_n\right)$ n'est pas monotone.&\textbf{d.~~}  Le sens de variation de la suite $\left(b_n\right)$ dépend de $b_0$.
%\end{tabularx}
%\end{center}
On a quel que soit le naturel $n$, \: $b_{n+1} - b_n = \ln \left(\dfrac{2}{\left(b_n \right)^2 + 3}\right)$.

Or quel que soit le réel $b_n$, \: $b_n^2 \geqslant 0$, donc $\left(b_n\right)^2 + 3 \geqslant 3 > 2$, donc $\dfrac{2}{\left(b_n \right)^2 + 3} < 1$ et enfin $\ln \left(\dfrac{2}{\left(b_n \right)^2 + 3}\right) < 0$.

Conclusion : $b_{n+1} - b_n < 0$ montre que la suite $\left(b_n\right)$ est décroissante.

\item %On considère la fonction $g$ définie sur l'intervalle $]0~;~+\infty[$ par : 

\[g(x) = \dfrac{\e^x}{x}.\]

%On note $\mathcal{C}_g$ la courbe représentative de la fonction $g$ dans un repère orthogonal.
%
%La courbe $\mathcal{C}_g$ admet :
%
%\begin{center}
%\begin{tabularx}{\linewidth}{X X}
%\textbf{a.~~}une asymptote verticale
%et une asymptote horizontale.&\textbf{b.~~}une asymptote verticale
%et aucune asymptote horizontale.\\
%\textbf{c.~~}aucune asymptote verticale et une asymptote horizontale.&\textbf{d.~~}aucune asymptote verticale et aucune asymptote horizontale.
%\end{tabularx}
%\end{center}
$\bullet~~$Au voisinage de zéro : $\displaystyle\lim_{x \to 0} \e^x = 1$ et $\displaystyle\lim_{x \to 0}x  = 0$, donc $\displaystyle\lim_{x \to 0} g(x) = +\infty$ : géométriquement l'axe des ordonnées est asymptote verticale à $\mathcal{C}_g$.

$\bullet~~$Au voisinage de plus l'infini : on sait que $\displaystyle\lim_{x \to + \infty} g(x) = + \infty$, donc pas d'asymptote.

\item %Soit $f$ la fonction définie sur $\R$ par

\[f(x) = x\e^{x^2+1}.\]

%Soit $F$ une primitive sur $\R$ de la fonction $f$. Pour tout réel $x$, on a :

%\begin{center}
%\begin{tabularx}{\linewidth}{X X}
%\textbf{a.~~}$F(x) = \dfrac12x^2\e^{x^2+1}$ 	& \textbf{b.~~}$F(x) = \left(1 + 2x^2 \right)\e^{x^2+1}$ \\
%\textbf{c.~~}$F(x) = \e^{x^2+1}$				& \textbf{d.~~}$F(x) = \dfrac12\e^{x^2+1}$
%\end{tabularx}
%\end{center}
En posant $u(x) = x^2 + 1$, on a $u'(x) = 2x$ et l'on peut écrire $f(x) = \dfrac12 \times 2x \times \e^{x^2+1} = \dfrac12 u'\e^u$ : on reconnait la dérivée de $\dfrac12\e^u = \dfrac12\e^{x^2 + 1}$.

Donc $F(x) = \dfrac12\e^{x^2+1}$.
\end{enumerate}
